\documentclass[lettersize,journal]{IEEEtran}

% -----------------
% Packages
% -----------------

\usepackage{amsmath,amsfonts}
\usepackage{algorithmic}
\usepackage{algorithm}
\usepackage{array}
\usepackage[caption=false,font=normalsize,labelfont=sf,textfont=sf]{subfig}
\usepackage{tabularx}
\usepackage{textcomp}
\usepackage{stfloats}
\usepackage{url}
\usepackage{verbatim}
\usepackage{graphicx}
\usepackage{cite}

% -----------------
% Global Defaults
% -----------------

\graphicspath{{figures/}}
% -----------------
% Custom Commands
% -----------------
\newcommand{\R}{\mathbb{R}}
\newcommand{\vect}[1]{\mathbf{#1}}
% -----------------
% Formatting
% -----------------

% \hyphenation{op-tical net-works semi-conduc-tor IEEE-Xplore}

\newcolumntype{Y}{>{\raggedright\arraybackslash}X}

\def\BibTeX{{\rm B\kern-.05em{\sc i\kern-.025em b}\kern-.08em
    T\kern-.1667em\lower.7ex\hbox{E}\kern-.125emX}}


% updated with editorial comments 8/9/2021

\begin{document}

\title{A Sample Article Using IEEEtran.cls\\ for IEEE Journals and Transactions}

\author{IEEE Publication Technology,~\IEEEmembership{Staff,~IEEE,}
        % <-this % stops a space
\thanks{This paper was produced by the IEEE Publication Technology Group. They are in Piscataway, NJ.}% <-this % stops a space
\thanks{Manuscript received April 19, 2021; revised August 16, 2021.}}

% The paper headers
\markboth{Journal of \LaTeX\ Class Files,~Vol.~14, No.~8, August~2021}%
{Shell \MakeLowercase{\textit{et al.}}: A Sample Article Using IEEEtran.cls for IEEE Journals}

\IEEEpubid{0000--0000/00\$00.00~\copyright~2021 IEEE}
% Remember, if you use this you must call \IEEEpubidadjcol in the second
% column for its text to clear the IEEEpubid mark.

\maketitle

\begin{abstract}
This document describes the most common article elements and how to use the IEEEtran class with \LaTeX \ to produce files that are suitable for submission to the IEEE.  IEEEtran can produce conference, journal, and technical note (correspondence) papers with a suitable choice of class options. 
\end{abstract}

\begin{IEEEkeywords}
Article submission, IEEE, IEEEtran, journal, \LaTeX, paper, template, typesetting.
\end{IEEEkeywords}

% \section{Introduction}

Recommender systems have evolved from early accuracy-focused algorithms 
to complex industrial systems across e-commerce, media, and social platforms.
Despite these advances, a tension exists between short-term business-driven goals and long-term user wellbeing.
Most current systems rely on easily measurable signals like clicks and watch time, 
which maximize immediate engagement but often do not explicitly model the user's cognitive state.

\subsection*{The Challenge of Engagement Bias}

Standard sequential and deep learning models treat user histories as ordered sequences to predict the next likely interaction.
While effective at capturing recency effects, these models implicitly assume that maximizing immediate clicks corresponds to maximizing satisfaction.
This approach neglects psychological factors such as \textbf{fatigue} and \textbf{diminishing utility}, 
where repeated exposure to similar content reduces its value and this leads to platform disengagement.

\subsection*{Proposed Solution: DS-HCRS}

To address these limitations, this work proposes a framework called \textbf{DS-HCRS} 
which goes beyond a single static strategy by incorporating:

\begin{itemize}
    \item \textbf{Dual-System Architecture:} Balancing a fast "impulsive" model for relevance with a reflective "long-term" optimizer.
    \item \textbf{Dynamics of Preference:} Estimating changes in user preference using recurrent models.
    \item \textbf{Cognitive State Modeling:} Estimating user fatigue and discontentment using Reinforcement Learning objectives.
\end{itemize}

These three systems alone allows the algorithm to implement more complex systems like 
updating user preference (embeddings) on the go, and cold start handling for new users and items.
By framing recommendation as a cognitive process rather than just a prediction problem,
we aim to develop systems that are sustainable, adaptive, and better aligned with long-term user experience.

\section{Introduction}

\emph{PLANNER}

\subsection{Background: Evolution of Recommender Systems}

Discuss progression:

\begin{enumerate}
    \item Collaborative Filtering
    \item Matrix Factorization
    \item Deep Learning RecSys
    \item Sequential Recommendation
    \item Reinforcement Learning RecSys
    \item LLM-Enhanced Recommenders
    \item Human-Centered Recommendation
\end{enumerate}

Goal: show historical trajectory toward your idea.

\subsection{Motivation}

Explain modern platform reality:

\begin{itemize}
    \item Engagement optimization dominates
    \item Click != satisfaction
    \item Platforms ignore cognitive limits
    \item Short-term optimization causes long-term degradation
\end{itemize}

Introduce concepts:

\begin{itemize}
    \item overconsumption
    \item user fatigue
    \item algorithmic addiction
    \item welfare-aware AI
\end{itemize}

\subsection{Problem Context}

Explain mismatch:

\begin{table}[!t]
\caption{An Example of a Table\label{tab:tableABDDC}}
\centering
\begin{tabular}{|l||c|}
\hline
Platform Objective & Human Objective \\
\hline
Maximize clicks & Maintain satisfaction \\
Short-term metrics & Long-term wellbeing \\
Immediate reward & Sustainable utility \\
\hline
\end{tabular}
\end{table}

\subsection{Research Gap Overview}

High-level gaps identified from surveyed papers:

\begin{itemize}
    \item Lack of fatigue modeling
    \item Static recommendation policies
    \item Weak long-term reward modeling
    \item No operational System-1/System-2 architecture
\end{itemize}
    
\subsection{Objectives}

Example:

\begin{itemize}
    \item Analyze limitations of engagement-centric recommenders    
    \item Study long-term optimization approaches
    \item Identify missing cognitive modeling components
    \item Propose DS-HCRS framework
\end{itemize}
    
\subsection{Contributions}

State clearly:

\begin{enumerate}
    \item Human-centric recommender formulation
    \item Fatigue-aware modeling framework
    \item Dual cognitive system architecture
    \item Long-term RL optimization perspective
\end{enumerate}

\section{Foundations and Background Concepts}

\section{Foundations and Background Concepts}

\textbf{PLANNER}

\subsection{Classical Recommendation Paradigms}

\begin{itemize}
    \item Content-based
    \item Collaborative filtering
    \item Matrix factorization
    \item Ranking objectives
\end{itemize}

\subsection{Sequential Recommendation}

Discuss:

\begin{itemize}
    \item GRU4Rec
    \item SASRec
    \item Transformer-based models
\end{itemize}

Key limitation:
$\Rightarrow$ assumes monotonic engagement utility.

\subsection{Reinforcement Learning in RecSys}

Explain:

\begin{itemize}
    \item MDP formulation
    \item reward shaping
    \item bandits vs RL
    \item long-term value estimation
\end{itemize}

Link to papers on long-term welfare optimization.

\subsection{Human Decision Theory (System 1 vs System 2)}

Introduce psychological grounding:

\begin{itemize}
    \item fast vs reflective cognition
    \item impulsive consumption vs reflective preference
\end{itemize}

Explain why RecSys implicitly models only System 1.

\subsection{Human-Centered AI \& Welfare Optimization}

Topics:

\begin{itemize}
    \item fairness
    \item wellbeing metrics
    \item ethical recommender design
    \item satisfaction vs engagement
\end{itemize}

\subsection{LLMs in Recommendation Systems}

Based on:

\begin{itemize}
    \item Awesome-LLM-for-RecSys repo
    \item recent arXiv papers
\end{itemize}

Discuss:

\begin{itemize}
    \item embeddings
    \item reasoning-based recommendation
    \item cold-start handling
\end{itemize}

% \section{Literature Survey}

\textbf{PLANNER}

This should NOT be paper-by-paper summaries.
Instead organize by research themes.

\subsection{Engagement-Optimized Recommender Systems}

Survey:

\begin{itemize}
    \item deep ranking models
    \item click prediction frameworks
\end{itemize}   

Findings:

\begin{itemize}
    \item short-term bias
    \item reward misalignment
\end{itemize}
    
\subsection{Sequential \& Session-Based Recommendation}

Focus:

\begin{itemize}
    \item temporal dynamics    
    \item user behavior modeling
\end{itemize}   

Limitation:

\begin{itemize}
    \item no cognitive saturation modeling.
\end{itemize}

\subsection{Reinforcement Learning for Long-Term Optimization}

Include:

\begin{itemize}
    \item contextual bandits    
    \item long-horizon RL papers
    \item welfare-aware RL
\end{itemize}
    
Findings:

\begin{itemize}
    \item reward definition still engagement-heavy.
\end{itemize}
    
\subsection{Human-Centered and Welfare-Aware Recommenders}

Discuss:

\begin{itemize}
    \item System-2 recommenders
    \item user satisfaction modeling
    \item ethical optimization
\end{itemize}
    
Key insight:  
$\Rightarrow$ theory exists, integration missing.

\subsection{Bias, Manipulation, and Overconsumption Studies}

From SemanticScholar papers:

\begin{itemize}
    \item addictive interaction dynamics
    \item divisive content amplification
    \item welfare tradeoffs    
\end{itemize}    

\subsection{LLM-Augmented Recommender Systems}

Themes:

\begin{itemize}
    \item semantic reasoning    
    \item concept embeddings
    \item controllable recommendation
\end{itemize}    

Gap:

\begin{itemize}
    \item reasoning used for relevance, not wellbeing.
\end{itemize}
    
\subsection{Key Trends from Major Conferences}

(Summarize trends rather than papers)

From:

\begin{itemize}
    \item RecSys 2024/2025
    \item SIGIR 2024
    \item KDD 2024
\end{itemize}
    

Observed shifts:

\begin{itemize}
    \item long-term evaluation
    \item controllability
    \item explainability
    \item hybrid AI systems
\end{itemize}

\subsection{Literature Synthesis Table (VERY IMPORTANT)}

Create table:

\begin{table}[!t]
\caption{An Example of a Table\label{tab:tableABC}}
\centering
\begin{tabular}{|l||c||c||c||c|}
\hline
Category & Representative Work & Objective & Strength & Limitation \\
\hline
Value 1 & Value 2 & Value 3 & Value 4 & Value 5 \\
\hline
\end{tabular}
\end{table}

This becomes your strongest survey artifact.

\section{Literature Survey}

\textbf{PLANNER}

This should NOT be paper-by-paper summaries.
Instead organize by research themes.

\subsection{Engagement-Optimized Recommender Systems}

Survey:

\begin{itemize}
    \item deep ranking models
    \item click prediction frameworks
\end{itemize}   

Findings:

\begin{itemize}
    \item short-term bias
    \item reward misalignment
\end{itemize}
    
\subsection{Sequential \& Session-Based Recommendation}

Focus:

\begin{itemize}
    \item temporal dynamics    
    \item user behavior modeling
\end{itemize}   

Limitation:

\begin{itemize}
    \item no cognitive saturation modeling.
\end{itemize}

\subsection{Reinforcement Learning for Long-Term Optimization}

Include:

\begin{itemize}
    \item contextual bandits    
    \item long-horizon RL papers
    \item welfare-aware RL
\end{itemize}
    
Findings:

\begin{itemize}
    \item reward definition still engagement-heavy.
\end{itemize}
    
\subsection{Human-Centered and Welfare-Aware Recommenders}

Discuss:

\begin{itemize}
    \item System-2 recommenders
    \item user satisfaction modeling
    \item ethical optimization
\end{itemize}
    
Key insight:  
$\Rightarrow$ theory exists, integration missing.

\subsection{Bias, Manipulation, and Overconsumption Studies}

From SemanticScholar papers:

\begin{itemize}
    \item addictive interaction dynamics
    \item divisive content amplification
    \item welfare tradeoffs    
\end{itemize}    

\subsection{LLM-Augmented Recommender Systems}

Themes:

\begin{itemize}
    \item semantic reasoning    
    \item concept embeddings
    \item controllable recommendation
\end{itemize}    

Gap:

\begin{itemize}
    \item reasoning used for relevance, not wellbeing.
\end{itemize}
    
\subsection{Key Trends from Major Conferences}

(Summarize trends rather than papers)

From:

\begin{itemize}
    \item RecSys 2024/2025
    \item SIGIR 2024
    \item KDD 2024
\end{itemize}
    

Observed shifts:

\begin{itemize}
    \item long-term evaluation
    \item controllability
    \item explainability
    \item hybrid AI systems
\end{itemize}

\subsection{Literature Synthesis Table (VERY IMPORTANT)}

Create table:

\begin{table}[!t]
\caption{An Example of a Table\label{tab:tableABC}}
\centering
\begin{tabular}{|l||c||c||c||c|}
\hline
Category & Representative Work & Objective & Strength & Limitation \\
\hline
Value 1 & Value 2 & Value 3 & Value 4 & Value 5 \\
\hline
\end{tabular}
\end{table}

This becomes your strongest survey artifact.

% \section{Problem Identification}

The literature survey highlights several limitations, a result of how these systems are designed and evaluated.
Current approaches prioritize short-term engagement and observable signals, often ignoring aspects 
like the human behavior (cognition, preferences, and satisfaction).
These limitations also lead to unintended consequences such as user fatigue, repetitive recommendations, 
and declining long-term engagement.

This section summarizes the key problems identified in the literature survey 
and motivates the need for a more human-centered recommendation framework.

%%%%%%%%%%%%%%%%%%%%%%%%%%%%%%%%%%%%%%%%%%%%%%%%%%%%%%%%%%%%%%%%
%%%%%%%%%%%%%%%%%%%%%%%%%%%%%%%%%%%%%%%%%%%%%%%%%%%%%%%%%%%%%%%%

\subsection{Short-Term Engagement Bias}

The focus on immediate user interaction has heavily influenced the design toward 
metrics like click-through rate (CTR), watch time, or the probability that a user
will interact with the next recommended item.
While these metrics are easy to measure and optimize for; they primarily represent 
short-term engagement rather than long-term satisfaction.

A recommendation that successfully triggers a click does not necessarily mean that
the user found it valuable or enjoyable; which in turn does not guarantee their 
retension over an extended period of time.
Additionally, users may interact with content out of curiosity rather than genuine interest.

Because of this, many recommendation algorithms are trained to 
maximize short-term interactions rather than long-term user value.
Over time, this can encourage content that attracts attention 
quickly but may not contribute satisfying user experiences.

This creates a bias in recommendation systems: \textbf{the system becomes optimized for immediate engagement rather than sustainable user satisfaction.}

%%%%%%%%%%%%%%%%%%%%%%%%%%%%%%%%%%%%%%%%%%%%%%%%%%%%%%%%%%%%%%%%
%%%%%%%%%%%%%%%%%%%%%%%%%%%%%%%%%%%%%%%%%%%%%%%%%%%%%%%%%%%%%%%%

\subsection{Lack of Cognitive State Modeling}

Traditional algorithms of recommendation work around 
available historical data about the user interactions.
In these models, a user is typically represented as 
a sequence of past clicks, views, or purchases.

This approach is effective for identifying behavioral patterns; however, it treats users 
as static data sources, completely ignoring their changing mental states and intentions.
Important psychological factors such as attention, curiosity, 
boredom, and fatigue does not have an explicit representation in the model.

As a result, the system can observe \textit{what users do}, but it 
cannot infer \textit{how users feel} about their current experience.
For example, a user who continues interacting with recommendations might appear 
highly engaged, even if they are becoming tired of the repetitive content.

Without modeling a cognitive state, recommender systems lack the ability to adapt 
to a strategy which can work around how users are experiencing the process.

%%%%%%%%%%%%%%%%%%%%%%%%%%%%%%%%%%%%%%%%%%%%%%%%%%%%%%%%%%%%%%%%
%%%%%%%%%%%%%%%%%%%%%%%%%%%%%%%%%%%%%%%%%%%%%%%%%%%%%%%%%%%%%%%%

\subsection{Absence of Fatigue and Diminishing Utility}

Another important limitation observed is the assumption that 
user utility remains constant across repeated interactions.
Most models assume that if a user interacted with some content in the past, 
recommending similar items repeatedly is a good move.

However, in reality, this repeatition often leads to \textit{diminishing utility}.
As users encounter similar recommendations multiple times, 
the perceived value of those recommendations decreases.
This leads to boredom, reduced interest, and eventually disengagement.

For example, if a user is recommended several videos from a single category or receives
similar product recommendations, the novelty of those recommendations gradually declines.
Despite that, many models continue to reinforce the same patterns 
because they rely heavily on past interaction signals.

Since this notion of fatigue is not explicitly modeled, recommender systems 
unintentionally generate overly repetitive recommendation streams.

%%%%%%%%%%%%%%%%%%%%%%%%%%%%%%%%%%%%%%%%%%%%%%%%%%%%%%%%%%%%%%%%
%%%%%%%%%%%%%%%%%%%%%%%%%%%%%%%%%%%%%%%%%%%%%%%%%%%%%%%%%%%%%%%%

\subsection{Static Recommendation Policies}

Most model from the traditional setting learn a single strategy representing the training data.
Therefore, given the historical pretext, these models perform really well.
Unfortunately, a user is a dynamic entity, with evolving flavours and tastes.
A single strategy is insufficient to keep track of a user's state or context.

This means the system is incapable of dynamically adjusting its decision strategy 
based on factors such as session length, user exploration patterns, or signs of fatigue.
Once a ranking model is trained, it is used to generate recommendations in the same way for every interaction.

In practice, different situations may require different recommendation strategies.
Early in a session, users may prefer familiar and relevant content.
Later in the session, introducing novelty or diversity might help prevent fatigue.

Since current systems typically rely on static ranking policies, they lack 
the ability to adapt its strategy to changing user cognitive state.

%%%%%%%%%%%%%%%%%%%%%%%%%%%%%%%%%%%%%%%%%%%%%%%%%%%%%%%%%%%%%%%%
%%%%%%%%%%%%%%%%%%%%%%%%%%%%%%%%%%%%%%%%%%%%%%%%%%%%%%%%%%%%%%%%

\subsection{Misaligned Evaluation Metrics}

Evaluation metrics play a critical role in shaping recommender system design.
Most systems are evaluated using metrics such as Precision@K, Recall@K, and 
NDCG; these metrics measure how accurately the system predicts user interactions.

These metrics are useful for benchmarking algorithm performance, however, 
only for predictive accuracy, rather than long-term user outcomes.
A system with a high ranking accuracy can still produce 
repetitive content, thus reducing user satisfaction.

Furthermore, offline evaluation methods often fail to capture the dynamic 
feedback loops that take place with real user interactions over time.
As a result, models that perform well in offline experiments 
may not always produce the best long-term user experiences.

%%%%%%%%%%%%%%%%%%%%%%%%%%%%%%%%%%%%%%%%%%%%%%%%%%%%%%%%%%%%%%%%
%%%%%%%%%%%%%%%%%%%%%%%%%%%%%%%%%%%%%%%%%%%%%%%%%%%%%%%%%%%%%%%%

\subsection{Summary of Identified Problems}

The analysis above highlights several key limitations in current recommender system design:

\begin{table}[h]
\caption{TODO Add a caption\label{tab:TODOlabel0102526}}
\centering
\begin{tabularx}{\columnwidth}{|l||Y|}
\hline
Problem & Description \\
\hline
Short-term engagement bias & Systems prioritize immediate interactions rather than long-term satisfaction \\
Lack of cognitive state modeling & Users are represented only through behavior sequences \\
Absence of fatigue modeling & Repeated recommendations may reduce user interest over time \\
Static recommendation policies & Systems do not adapt strategies based on user state \\
Misaligned evaluation metrics & Metrics reward interaction prediction rather than user wellbeing \\
\hline
\end{tabularx}
\end{table}

Together, these limitations suggest that current recommender systems are primarily optimized for \textbf{short-term behavioral signals rather than long-term human-centered outcomes}.

These challenges motivate the development of new recommendation frameworks that can explicitly consider user cognitive state, fatigue, and long-term satisfaction when generating recommendations.


\section{Problem Identification}

\textbf{PLANNER}

This section converts survey → research need.

\subsection{Structural Problems in Current Recommenders}

\subsubsection{Short-Term Reward Bias}

Explain mathematically.

\subsubsection{Absence of Cognitive State Modeling}

Users treated as static preference vectors.

\subsubsection{No Fatigue or Saturation Representation}

Utility assumed additive.

\subsubsection{Static Decision Policies}

No adaptive reasoning depth.

\subsubsection{Misalignment Between Platform Metrics and User Welfare}

\subsection{Empirical Consequences}

\begin{itemize}
    \item binge consumption    
    \item reduced novelty
    \item churn
    \item polarization
\end{itemize}
    
\subsection{Formal Problem Definition}

Define:

User state:

\begin{equation}
    s_t = (history, context, fatigue)
\end{equation}

Objective:

\begin{equation}
    \max \mathbb{E} \sum_{t} R_{t}^{long-term}  
\end{equation}

Explain why classical models fail.

% \section{Proposed Methodology: DS-HCRS Framework}

\textbf{PLANNER}

(Conceptual - no implementation details.)

\subsection{Design Philosophy}

Principles:

\begin{itemize}
    \item human-aware optimization
    \item adaptive cognition modeling
    \item long-term value
\end{itemize}
    
\subsection{System Overview Architecture}

High-level diagram description:

User → State Modeling → Dual Systems → Gating → Recommendation → Feedback.

\subsection{System 1: Fast Engagement Model}

Role:

\begin{itemize}
    \item reactive ranking
    \item immediate reward modeling.
\end{itemize}
    
\subsection{Fatigue \& Cognitive Burden Model}

Define variables:

\begin{itemize}
    \item repetition
    \item entropy decay
    \item session length
    \item novelty loss
\end{itemize}    

Explain theoretical grounding.

\subsection{System 2: Reflective Long-Term Optimizer}

Conceptual RL formulation.

Explain:

\begin{itemize}
    \item delayed reward
    \item retention proxy
    \item welfare regularization.
\end{itemize}    

\subsection{Dual-System Gating Mechanism}

Decision policy rationale.

Explain adaptive reasoning depth idea.

\subsection{Optional Concept Embedding Module}

LLM embeddings for cold start.

\subsection{Learning Objective Formulation}

General reward:

\[
R_t = \alpha E_t + \beta L_t - \gamma F_t
\]

Explain meaning conceptually.

\subsection{Expected Advantages Over Existing Systems}

Table comparison:

\begin{table}[!t]
\caption{An Example of a Table\label{tab:table1}}
\centering
\begin{tabular}{|l||c||c|}
\hline
Feature & Traditional & Proposed \\
\hline
Fatigue modeling & cross & tick \\
Long-term reward & Partial & Explicit \\
Adaptive reasoning & cross & tick \\
\hline
\end{tabular}
\end{table}
\section{Proposed Methodology: DS-HCRS Framework}

\textbf{PLANNER}

(Conceptual - no implementation details.)

\subsection{Design Philosophy}

Principles:

\begin{itemize}
    \item human-aware optimization
    \item adaptive cognition modeling
    \item long-term value
\end{itemize}
    
\subsection{System Overview Architecture}

High-level diagram description:

User -> State Modeling -> Dual Systems -> Gating -> Recommendation -> Feedback.

\subsection{System 1: Fast Engagement Model}

Role:

\begin{itemize}
    \item reactive ranking
    \item immediate reward modeling.
\end{itemize}
    
\subsection{Fatigue \& Cognitive Burden Model}

Define variables:

\begin{itemize}
    \item repetition
    \item entropy decay
    \item session length
    \item novelty loss
\end{itemize}    

Explain theoretical grounding.

\subsection{System 2: Reflective Long-Term Optimizer}

Conceptual RL formulation.

Explain:

\begin{itemize}
    \item delayed reward
    \item retention proxy
    \item welfare regularization.
\end{itemize}    

\subsection{Dual-System Gating Mechanism}

Decision policy rationale.

Explain adaptive reasoning depth idea.

\subsection{Optional Concept Embedding Module}

LLM embeddings for cold start.

\subsection{Learning Objective Formulation}

General reward:

\begin{equation}
    R_t = \alpha E_t + \beta L_t - \gamma F_t
\end{equation}

Explain meaning conceptually.

\subsection{Expected Advantages Over Existing Systems}

Table comparison:

\begin{table}[!t]
\caption{An Example of a Table\label{tab:table1}}
\centering
\begin{tabular}{|l||c||c|}
\hline
Feature & Traditional & Proposed \\
\hline
Fatigue modeling & cross & tick \\
Long-term reward & Partial & Explicit \\
Adaptive reasoning & cross & tick \\
\hline
\end{tabular}
\end{table}

% \section{Evaluation Design}

\textbf{PLANNER}

(No experiments yet.)

\subsection{Baseline Models}

\begin{itemize}
    \item sequential recommender
    \item RL recommender
\end{itemize}

\subsection{Metrics}

\subsubsection{Short-term}

\begin{itemize}
    \item NDCG
    \item Recall@K
\end{itemize}

\subsubsection{Long-term}

\begin{itemize}
    \item retention proxy
    \item churn delay
    \item fatigue score
\end{itemize}

\subsection{Simulation Framework}

Explain trajectory simulation conceptually.

\subsection{Ablation Philosophy}

Why component removal validates design.


\section{Evaluation Design}

\textbf{PLANNER}

(No experiments yet.)

\subsection{Baseline Models}

\begin{itemize}
    \item sequential recommender
    \item RL recommender
\end{itemize}

\subsection{Metrics}

\subsubsection{Short-term}

\begin{itemize}
    \item NDCG
    \item Recall@K
\end{itemize}

\subsubsection{Long-term}

\begin{itemize}
    \item retention proxy
    \item churn delay
    \item fatigue score
\end{itemize}

\subsection{Simulation Framework}

Explain trajectory simulation conceptually.

\subsection{Ablation Philosophy}

Why component removal validates design.

% \section{Expected Research Contributions}

EVERYTHING IS A I HERE, NEEDS REWRITING

The proposed Dual-System Human-Centric Recommender System (DS-HCRS) introduces several conceptual and methodological contributions to the design of recommender systems.
While the implementation and empirical evaluation are planned for future work, the proposed framework outlines a structured approach for incorporating cognitive modeling and long-term optimization into recommendation pipelines.

The expected contributions of this work can be summarized as follows.

\subsection{Dual-System Recommendation Architecture}

This work proposes a dual-system recommender architecture inspired by cognitive theories of human decision-making.
The framework separates recommendation strategies into two complementary components:

\begin{itemize}
    \item \textbf{System 1}, which captures short-term behavioral preferences and predicts immediate user interactions.
    \item \textbf{System 2}, which evaluates recommendations using longer-term utility signals and broader contextual information.
\end{itemize}

By explicitly separating impulsive interaction prediction from reflective optimization, the proposed system introduces a structured way to balance short-term engagement with long-term user outcomes.
This architectural design provides a flexible framework that can integrate different recommendation models while maintaining a clear conceptual separation of objectives.

\subsection{Fatigue-Aware Recommendation Modeling}

Another key contribution of this work is the introduction of an explicit \textit{user fatigue and cognitive state model} within the recommendation pipeline.

Traditional recommender systems assume that additional user interactions always indicate positive engagement.
However, prolonged sessions and repeated exposure to similar content may lead to boredom or cognitive overload.
The proposed fatigue modeling component attempts to estimate this state using observable interaction signals such as repetition, novelty, and session length.

By incorporating fatigue estimates into the recommendation process, the system can adapt its behavior when signs of declining user interest are detected.

\subsection{Adaptive Strategy Selection Through Gating}

The proposed system introduces an adaptive gating mechanism that dynamically determines how recommendations from different components should be combined.

Instead of relying on a single recommendation strategy for all situations, the system adjusts the influence of System 1 and System 2 based on the user's current interaction context.
This adaptive decision process allows the recommender to maintain efficiency during normal usage while introducing more reflective reasoning when the user's experience begins to deteriorate.

Such adaptive control mechanisms provide a potential pathway toward more context-aware recommender systems.

\subsection{Integration of Long-Term Utility Optimization}

The proposed framework incorporates long-term optimization principles into the recommendation pipeline by modeling recommendations as sequential decisions that influence future user behavior.

By combining short-term engagement signals with longer-term indicators such as retention and fatigue, the framework outlines a multi-objective reward structure that can guide recommendation policies toward more sustainable outcomes.

This integration highlights how reinforcement learning-based ideas can be used to move beyond purely engagement-driven recommendation objectives.

\subsection{Human-Centered Evaluation Perspective}

Finally, this work proposes an evaluation framework that extends beyond traditional recommendation accuracy metrics.
In addition to standard ranking measures, the evaluation strategy includes metrics related to diversity, novelty, fatigue indicators, and long-term engagement.

This broader evaluation perspective aligns the system with recent research directions in human-centered recommender systems and responsible AI.
By incorporating metrics related to user experience and cognitive load, the proposed evaluation framework aims to provide a more comprehensive understanding of recommender system behavior.

Overall, the proposed framework contributes a structured approach for combining behavioral prediction, cognitive modeling, and long-term optimization within a single recommender system architecture.
The design emphasizes adaptability and human-centered objectives while remaining compatible with existing recommendation techniques.



\section{Expected Research Contributions}

\textbf{PLANNER}

Conceptual contributions:

\begin{itemize}
    \item Cognitive-aware recommendation formulation
    \item Fatigue-aware utility modeling
    \item Dual-system control paradigm
    \item Human-aligned evaluation framing
\end{itemize}

% \section{Limitations and Ethical Considerations}

\textbf{PLANNER}

Very important academically.

Discuss:

\begin{itemize}
    \item proxy wellbeing limitations
    \item offline evaluation constraints
    \item simulation assumptions
    \item ethical deployment risks
\end{itemize}

\section{Limitations and Ethical Considerations}

\textbf{PLANNER}

Very important academically.

Discuss:

\begin{itemize}
    \item proxy wellbeing limitations
    \item offline evaluation constraints
    \item simulation assumptions
    \item ethical deployment risks
\end{itemize}

% \section{Future Research Directions}

\textbf{PLANNER}

Examples:

\begin{itemize}
    \item real user studies
    \item adaptive psychological modeling
    \item causal preference learning
    \item multi-user ecosystem welfare
\end{itemize}

\section{Future Research Directions}

\textbf{PLANNER}

Examples:

\begin{itemize}
    \item real user studies
    \item adaptive psychological modeling
    \item causal preference learning
    \item multi-user ecosystem welfare
\end{itemize}

% \section{Conclusion}

\textbf{PLANNER}

Summarize:

Problem → Gap → DS-HCRS → Expected impact.

\section{Conclusion}

\textbf{PLANNER}

Summarize:

Problem → Gap → DS-HCRS → Expected impact.

\bibliographystyle{IEEEtran}
\bibliography{references}

\vfill

\end{document}


